Die Durchschnittsgeschwindigkeit ist die gesamte Strecke durch die gesamte Zeit. Die beiden Teilstrecken dauern
\begin{align}
 t_a &= \frac{a}{u}, \qquad t_b = \frac{b}{v}
\end{align}
Damit ist
\begin{align}
 \bar v &= \frac{a+b}{t_a+t_b} = \frac{a+b}{\frac{a}{u}+\frac{b}{v}} = \result{\frac{(a+b)\,u\,v}{a\,v+b\,u}}
\end{align}
Für gleich lange Teilstrecken ($a=b$) ergibt sich $\bar v = \frac{2\,u\,v}{u+v}$. Das ist kleiner als der Mittelwert $\frac{u+v}{2}$ der beiden Geschwindigkeiten (ausser für $u=v$), weil man für die langsame Teilstrecke länger braucht.
